% !TEX root = main.tex
\subsection{Data acquisition and provenance}\label{sec:methods-data-acquisition}

A comparison between learning methods is only as interpretable as the data on
which it is made. Promoter benchmarks are distributed as files in code
repositories, journal supplements, and institutional servers. These copies can
change without notice, and several of them differ from the publications that
describe them. The data used in this thesis are therefore treated as an object
whose origin must be demonstrable, following the findability, accessibility,
interoperability and reusability principles for scientific data
\citep{wilkinson2016fair}. Three requirements guided the procedure: each
experimental example must be traceable to an identified upstream file; the
upstream copy must be identified independently of the moment at which it was
obtained; and any later acquisition must be distinguishable from the first.

\paragraph{Reference collection.}
The datasets released with the evaluated studies
\citep{umarov2017recognition,bhandari2021comparison,martinez2024cdbprom,zhang2024graphpro},
together with primary references used for reconstruction, such as the
Eukaryotic Promoter Database \citep{meylan2020epd}, were acquired again directly
from their distribution points into a single versioned reference collection.
Each source was identified by an access identifier that cannot drift. Code
repositories were pinned to a full commit identifier, and every file had to
reproduce the content hash that the repository records for it. For servers
without version control, the release number was recorded together with the
digests published by the server or recorded by earlier project contracts. Two
dates are kept apart: the date on which the project obtained each file, and the
date at which the source reports that its copy was last changed, such as the
commit date or the last-modification date declared by the server. For example,
the CNNProm benchmark obtained in 2026 corresponds to a repository state from
2016. Any change of source, version, or expected content defines a new
collection version rather than modifying the existing one.

\paragraph{Canonical representation.}
The acquired files were converted into one uniform tabular representation per
released dataset. The conversion was deliberately conservative. It removed
formatting whitespace and normalized letter case, and it neither discarded,
trimmed, deduplicated, nor relabelled any record. Properties that later stages
must handle, namely length deviations, symbols outside $\Sigma$, exact
duplicates, and identical sequences carrying opposite labels, were recorded as
attributes of each record instead of being corrected silently. A released
dataset was accepted only when its class counts equalled those of the release
and, where a previously used copy existed, when both copies contained the same
multiset of sequences in each class.

\paragraph{Continuity with earlier experiments.}
The experiments reported in this thesis were prepared from copies obtained
before this procedure existed, for which no acquisition record survives. Their
acquisition date is therefore reported as unknown, and it is not inferred from
file-system metadata. Their origin was instead established by evidence. The
preparation pipeline was executed again on the newly acquired files, under the
original rules and random seed. It reproduced, bit for bit, the model-ready
training and test partitions of all eighteen confirmatory datasets, and of
seven of the eight GraphPro datasets, which are retained only for exploratory
use. In the remaining GraphPro dataset, a correction of the organism name
changed the seed from which its partition is derived; the same examples were
recovered, assigned differently between partitions. The confirmatory results
are thus linked to identified upstream releases without being recomputed.

\paragraph{Observed properties of the releases.}
The renewed acquisition exposed release characteristics that are relevant for
interpreting published benchmarks. In one multi-species prokaryotic benchmark,
identical sequences occur with both labels, and some of them span the
distributed training and test partitions. In two eukaryotic promoter
collections, a large fraction of the positive examples are exact duplicates.
One file distributed as a plant promoter set carries mouse genome coordinates
in its headers, which suggests, but does not establish, a mislabelled release.
For one bacterial resource, the published window around the transcription
start site does not match the length of the distributed sequences. These
observations are reported as properties of the releases. They motivate the
explicit treatment of duplicates and label collisions summarized in
Table~\ref{tab:sequence-integrity-policies}, and they qualify any comparison with
results reported on the original files.

\paragraph{Limitations.}
For servers that publish no checksums, the identity of a file rests on its first
acquisition by the project, supplemented by agreement with earlier copies where
these exist. Several servers report no modification date, so the version of
their content can be dated only by acquisition. Exact agreement of sequences
does not establish their biological annotation: the procedure guarantees that
the data are the released data, not that the released labels are correct.
